% Appendix: LHB Evaluation Runtime Configuration
% Describes the parameters accepted when invoking `lhb run`.

\section{Loop Engineering Evaluation Runtime Configuration}
\label{app:run-config}
The \lhb{} evaluation runtime is invoked via \texttt{lhb run --config \textit{config.yaml}}, and the resulting run artifact records the loop parameters used for attribution.
The evaluation runtime treats the evaluation adapter, loop budget, and residual handoff policy as explicit loop engineering factors for each run.
All parameters can be set in the YAML configuration file or overridden by CLI flags. Table~\ref{tab:run-params} summarizes the runtime parameters used to separate attribution across model choice, evaluation adapter, budget, and residual handoff policy while keeping the evaluated task contract fixed.

\begin{table}[!htb]
  \centering
  \caption{Runtime parameters for \texttt{lhb run}.}
  \label{tab:run-params}
  \vspace{0.3em}
  \footnotesize
  \setlength{\tabcolsep}{3pt}
  \renewcommand{\arraystretch}{1.15}
  \resizebox{\columnwidth}{!}{%
  \begin{tabular}{@{}l l l p{6.5cm}@{}}
    \toprule
    \textbf{Group} & \textbf{Parameter} & \textbf{Default} & \textbf{Description} \\
    \midrule
    \multirow{3}{*}{Dataset}
      & \texttt{dataset\_path} & \texttt{tasks/} & Path to the task directory tree. \\
      & \texttt{task\_ids} & all               & Task ID globs to include (repeatable). \\
      & \texttt{exclude\_task\_ids} & \texttt{[]} & Task ID globs to exclude (repeatable). \\
    \midrule
    \multirow{4}{*}{Evaluation adapter}
      & \texttt{agent} & \texttt{oracle} & Built in evaluation adapter name. \\
      & \texttt{agent\_import\_path} & --             & Import path for a custom evaluation adapter class. \\
      & \texttt{model} & --                & LLM in \texttt{provider/model} format. \\
      & \texttt{agent\_kwargs} & \texttt{\{\}} & Key value pairs forwarded to the evaluation adapter. \\
    \midrule
    \multirow{3}{*}{Output}
      & \texttt{output\_path} & \texttt{runs/} & Root directory for run artifacts. \\
      & \texttt{run\_id} & timestamp         & Unique identifier for this run. \\
      & \texttt{n\_tasks} & all               & Cap on number of tasks to evaluate. \\
    \midrule
    \multirow{2}{*}{Build}
      & \texttt{no\_rebuild} & \texttt{false} & Skip rebuilding Docker images. \\
      & \texttt{cleanup} & \texttt{true} & Remove images after run completes. \\
    \midrule
    \multirow{2}{*}{Logging}
      & \texttt{log\_level} & \texttt{info} & Python logging level. \\
      & \texttt{livestream} & \texttt{false} & Stream evaluated loop I/O in real time. \\
    \midrule
    \multirow{2}{*}{Concurrency}
      & \texttt{n\_concurrent} & \texttt{4} & Tasks evaluated in parallel. \\
      & \texttt{n\_attempts} & \texttt{1} & Attempts per task. \\
    \midrule
    \multirow{3}{*}{Timeout}
      & \texttt{global\_timeout\_multiplier} & \texttt{1.0} & Multiplier for all per task timeouts. \\
      & \texttt{global\_agent\_timeout\_sec} & per task     & Override loop budget (seconds). \\
      & \texttt{global\_test\_timeout\_sec} & per task     & Override test budget in seconds. \\
    \bottomrule
  \end{tabular}}
\end{table}

\paragraph{Per task timeouts.}
Each task declares \texttt{max\_agent\_timeout\_sec} (default~1\,800\,s) and \texttt{max\_test\_timeout\_sec} (default~3\,600\,s) in its \texttt{task.yaml}.
The global overrides and multiplier allow uniform budget scaling without rewriting individual task files. Budget changes therefore remain attributable to run configuration rather than task materialization.

\paragraph{Residual handoff structure.}
The evaluation runtime executes an \emph{outer loop} over evaluation adapter runs. Each round starts a fresh shell session while preserving the workspace from prior rounds, so residual handoff is measured as an observed loop engineering factor rather than treated as an implicit retry mechanism.
The number of rounds is controlled by \texttt{outer\_loop\_count} (default~3); setting it to \texttt{0} or \texttt{unlimited} removes the cap, allowing repeated residual handoffs until the timeout budget is exhausted and making continuation policy explicit in the loop trace.
Within each round, rate limit errors trigger retries governed by \texttt{rate\_limit\_retries\_per\_round} (default~5) with exponential backoff starting at \texttt{rate\_limit\_backoff\_sec} (default~45\,s), recording provider availability separately from loop progress.

\paragraph{Model and provider configuration.}
Provider credentials and environment variables are supplied via a separate YAML file (\texttt{model\_config\_path} or \texttt{-{}-model-config}), which supports \texttt{env\_files} (dotenv paths) and an explicit \texttt{env} map. This separation places provider transport outside the loop engineering factors attributed to model choice and evaluation adapter.
Keys already set in the process environment are preserved unless \texttt{override\_process\_env:\allowbreak{} true} is selected, reducing cross run configuration drift.

\paragraph{Example configuration.}
The following minimal configuration specifies the dataset and execution policy while recording the evaluation adapter, loop budget, and residual handoff policy as explicit run metadata.

\begin{promptbox}{Run Configuration \textnormal{(\texttt{lhb.run.yaml})}}
dataset_path: tasks
output_path: runs

agent: mini-swe-agent
model: anthropic/claude-sonnet-4-20250514

n_concurrent: 4
n_attempts: 1

no_rebuild: false
cleanup: true
log_level: info
livestream: false

agent_kwargs:
  outer_loop_count: 3
  rate_limit_retries_per_round: 5
  rate_limit_backoff_sec: 45

global_timeout_multiplier: 1.0
\end{promptbox}

\begin{promptbox}{Model Environment \textnormal{(\texttt{lhb.model.yaml})}}
env_files: []
override_process_env: false
env:
  ANTHROPIC_API_KEY: "sk-ant-..."
\end{promptbox}
